% Generated draft: requires mathematical and source review.
\documentclass[11pt]{article}
\usepackage[margin=25mm]{geometry}
\usepackage{fontspec}
\setmainfont{Latin Modern Roman}
\usepackage{amsmath,amssymb,mathtools,unicode-math}
\setmathfont{Latin Modern Math}
\usepackage{enumitem,needspace,xurl,hyperref}
\setlength{\parindent}{0pt}
\setlength{\parskip}{5pt}
\setlength{\emergencystretch}{3em}
\newcommand{\N}{\mathbb{N}}
\newcommand{\Z}{\mathbb{Z}}
\newcommand{\Q}{\mathbb{Q}}
\newcommand{\R}{\mathbb{R}}
\newcommand{\C}{\mathbb{C}}
\newcommand{\F}{\mathbb{F}}
\newcommand{\A}{\mathbb{A}}
\newcommand{\PP}{\mathbb{P}}
\newcommand{\E}{\mathbb{E}}
\newcommand{\eps}{\varepsilon}
\newcommand{\dd}{\,\mathrm{d}}
\begin{document}
\section*{UnsolvedMath ID: 30002091}
\textbf{Source problem code:} catalogue entry
\subsection{Definitions and mathematical statement}
Write \(\pi:\mathcal X\to D\) for the deformation in the supplied record, where \(0\in D\) is its distinguished parameter, and set \(D^*=D\setminus\{0\}\) and \(X_t=\pi^{-1}(t)\). The source denotes exceptional curves in \(X_t\), for \(t\in D^*\), by \(E_{i,t}\), and states that \(X_0\) contains no exceptional curves. The index range and the manner in which the curves are followed as \(t\) varies are not supplied. The title places the question in the setting of class VII complex surfaces, but the precise preceding construction and hypotheses on the deformation are missing.

In the usual smooth complex surface terminology, an exceptional curve of the first kind is a smooth curve \(E\cong\mathbb P^1\) with self-intersection \(E^2=-1\). Whether this is exactly the intended meaning of ``exceptional curve'' in the extracted question requires source verification; it is not an additional asserted hypothesis here.

\textbf{Qualitative problem.} For the deformation and exceptional curves arising from the omitted source construction, describe the behavior of \(E_{i,t}\) as \(t\to0\), and explain the geometric mechanism by which exceptional curves present in the noncentral fibers fail to occur in \(X_0\). The intended notion of limiting behavior is not specified, so this is a qualified reconstruction of a question rather than a fully specified convergence assertion.

The source reports that the areas diverge to infinity in its example. If area is measured using Hermitian metrics with associated real \((1,1)\)-forms \(\omega_t\), this report would be expressed, for each intended tracked curve, as
\[
\operatorname{Area}_{\omega_t}(E_{i,t})=\int_{E_{i,t}}\omega_t\longrightarrow+\infty\qquad(t\to0).
\]
The example, metrics, normalization, parameter domain, and quantification over curve labels are not supplied. Consequently, this displayed formula only explicates the reported example-dependent behavior; it is not a universal conclusion or an independently verified assertion.
\subsection{Short English statement}
The supplied question asks how exceptional curves in the noncentral fibers of a class VII surface deformation behave when the central fiber has no exceptional curves. It reports area divergence in a particular example. The omitted construction, meaning of exceptional curve, choice of area, and notion of limit must be recovered before a fully precise mathematical formulation is possible.
\subsection{Sources}
\begin{itemize}
\item Supplied local UnsolvedMath registry record, permanent ID \(30002091\), exact source problem code \texttt{catalogue entry}, title ``Vanishing Exceptional Curves in Class-VII Surface Degenerations,'' in the Oberwolfach Reports Open Problems collection. The record gives proposed year \(2012\) and names no proposer. Its statement status is \texttt{reconstructed\_unverified}. The supplied verification note says that the official report confirms the qualitative degeneration question, but that its notation depends on a preceding construction and lacks sufficient local definitions for a self-contained statement. This verification claim is supplied metadata, not independent verification here.
\item Supplied report citation: \emph{Invariants in Low-Dimensional Topology and Knot Theory} (2013), Oberwolfach Reports, pp.~1687--1758, DOI \texttt{10.4171/owr/2012/28}. \url{https://doi.org/10.4171/owr/2012/28}. This is also the supplied literature link labeled ``original OWR report.'' The background describes Oberwolfach Reports as documenting mathematical workshops and their open-problem discussions.
\item Supplied literature reference: Teleman, \emph{Instantons and curves on class VII surfaces}. \url{https://arxiv.org/abs/0704.2634}.
\item Supplied dated literature curation, checked \texttt{2026-08-21}: the record labels the problem ``partially solved.'' Its assessment states that Donaldson and Teleman methods establish substantial class VII results, including curves and classification information for small second Betti number, while the extracted limiting behavior of the exceptional curves as the parameter tends to zero remains unresolved in full generality. This is supplied dated triage derived from the OWR curation, not an independently verified literature assessment; its time-sensitive status and cited claims require independent checking.
\end{itemize}
\end{document}
